\documentclass[11pt,a4paper]{article}
\usepackage[T1]{fontenc}
\usepackage[margin=25mm]{geometry}
\usepackage{amsmath,amssymb,amsthm}
\usepackage{microtype}

\newtheorem{conjecture}{Conjecture}
\newtheorem{proposition}{Proposition}
\theoremstyle{remark}
\newtheorem*{remark}{Remark}
\pagestyle{empty}
\setlength{\parindent}{0pt}
\setlength{\parskip}{0.45em}

\title{The Fair Traversal Conjecture}
\author{Luc Gillibert, GPT 6 Astra}
\date{17/11/2026}

\begin{document}
\maketitle
\vspace{-2em}

All graphs are finite, simple, and undirected. A graph is
\emph{vertex-transitive} if its automorphism group acts transitively on
its vertices. A \emph{Hamiltonian path} visits every vertex exactly
once; the one-vertex path is allowed.

\begin{conjecture}[Fair Traversal]
Every connected vertex-transitive graph with at least two vertices admits
a nonempty finite multiset of Hamiltonian paths in which every edge
occurs equally often.
\end{conjecture}

More explicitly, for every such graph $G=(V,E)$ there exist positive
integers $r,k$ and Hamiltonian paths $P_1,\ldots,P_r$, not necessarily
distinct, such that
\[
  \bigl|\{i\in\{1,\ldots,r\}:e\in E(P_i)\}\bigr|=k
  \qquad\text{for every }e\in E.
\]
Thus complete traversals can be chosen so that no edge is favored.
Writing $n=|V|$ and $m=|E|$, counting edge occurrences gives
$r(n-1)=km$. In particular, choosing an index uniformly from
$\{1,\ldots,r\}$ uses each edge with probability $(n-1)/m$.

\begin{conjecture}[Lov\'asz, Hamiltonian-path formulation]
Every connected vertex-transitive graph has a Hamiltonian path.
\end{conjecture}

\begin{proposition}
The Fair Traversal Conjecture implies Lov\'asz's conjecture.
\end{proposition}

\begin{proof}
Assume the Fair Traversal Conjecture, and let $G$ be a connected
vertex-transitive graph. If $G$ has one vertex, that vertex alone is a
Hamiltonian path. Otherwise, the Fair Traversal Conjecture supplies a
nonempty multiset $\{P_1,\ldots,P_r\}$ of Hamiltonian paths in $G$.
Since $r\geq 1$, the path $P_1$ exists and is a Hamiltonian path in $G$.
This proves Lov\'asz's conjecture.
\end{proof}

\begin{remark}
The implication uses only the nonemptiness of the collection. The
additional demand in Fair Traversal is equal edge usage. For example,
on a cycle with $n\geq 3$ vertices, deleting each edge in turn produces
$n$ Hamiltonian paths, and every edge belongs to exactly $n-1$ of them.
\end{remark}

\end{document}
